\section{The construction of the \texorpdfstring{\SLR{}}{SL(2,R)}/SO(2) coset}
\label{sec:sl2r-coset}

We briefly review the construction of the $\frac{\SLR{}}{SO(2)}$ coset representative, which is important for the discussion of the \SLR{} symmetry and the interplay between \SLZ{} and U(1)-violation.

The coset representative of $\frac{\SLR{}}{SO(2)}$ is an \SLR{} matrix $\mathcal{V}$ with a local SO(2) symmetry.
Generically, $\dim \SLR{} = 3$, so $\mathcal{V}$ is parameterised by three real scalars: $\mathcal{V} = \mathcal{V}(\tau_1, \tau_2, \theta)$.
It is convenient to separate the SO(2) factor, parameterised by $\theta$, using the Iwasawa decomposition \cite{iwasawa1949}
\begin{equation}
\mathcal{V} =
\begin{pmatrix}
1 & 0 \\
\tau_1 & 1
\end{pmatrix}
\begin{pmatrix}
\frac{1}{\sqrt{\tau_2}} & 0\\
0 & \sqrt{\tau_2}
\end{pmatrix}
\begin{pmatrix}
\cos\theta & \sin\theta\\
-\sin\theta & \cos\theta
\end{pmatrix}.
\label{Iwasawa-decomposition}
\end{equation}
In this representation, \SLR{} acts on $\mathcal{V}$ from the left, with $\Lambda \equiv \begin{pmatrix}
d&c\\b&a\end{pmatrix} \in \SLR{}$.
The representative $\mathcal{V}$ by itself is not physical due to the SO(2) gauge redundancy.
However, one can construct $\mathcal{V}\mathcal{V}^T$, which transforms covariantly under \SLR{}, removes the SO(2) factor, and gives the physical metric for the scalar manifold, parameterised by $\tau_1, \tau_2$:
\begin{equation}
\mathcal{M} = \mathcal{V}\mathcal{V}^T = \frac{1}{\tau_2}\begin{pmatrix}
1 & \tau_1\\
\tau_1 & |\tau|^2
\end{pmatrix}.
\end{equation}

To construct the kinetic term for the scalars, one builds the sl(2, R)-valued Maurer-Cartan 1-form $\mathcal{V}^{-1}\partial_m\mathcal{V}$, which decomposes into the non-compact and compact elements of the sl(2, R) algebra: $\mathcal{V}^{-1}\partial_m\mathcal{V} = P_m + Q_m$, where $P_m$ projects onto the non-compact generators and $Q_m$ is the SO(2) connection.
The kinetic term for the scalars is then given by $\mathrm{tr}(P^m P_m)$, which is positive definite, as ensured by tracing over the non-compact elements of a Lie algebra.
Computing $P, Q$, with $SO(2)\cong U(1)$ complexified via $e^{i\theta} = \cos\theta + i\sin\theta$, one finds
\begin{equation}
P = P_1\begin{pmatrix}
    0&1\\1&0
\end{pmatrix}+P_2\begin{pmatrix}
    1&0\\0&-1
\end{pmatrix},\quad
P_1 + iP_2 = e^{2i\theta}\frac{\partial \tau}{2 \tau_2},\quad
Q = \partial \theta - \frac{1}{2}\frac{\partial \tau_1}{\tau_2}.
\end{equation}
So the vielbein $P$ transforms covariantly under the local U(1) with charge 2, and the kinetic term $\frac{|\partial \tau|^2}{2\tau_2^2}$ is U(1) invariant.
This U(1) is independent of the global \SLR{} transformation, and is a gauge redundancy reflecting the choice of the U(1) frame.

The parameter $\theta$ here is the non-dynamical U(1) phase.
We can remove its presence in the action by picking a gauge $\theta = 0$, leaving two real scalars $\tau_1, \tau_2$ that parameterise the 2d coset $\frac{\SLR{}}{U(1)}$ without redundancy.
However, general \SLR{} transformations do not preserve the $\theta = 0$ gauge, so a compensating U(1) transformation given by $\theta = -\arg(c\tau+d)$ is induced when $\Lambda$ acts on $\mathcal{V}$ from the left.
This is independent of the local U(1) symmetry, which may be violated, but is part of the global \SLZ{} symmetry, which should always be preserved.

This is another way to view the breaking of the U(1) symmetry: the original U(1) is the SO(2) factor in \SLR{}, so when \SLR{} is broken to \SLZ{}, the U(1) is also broken.

Besides $P$, the other bosonic field that transforms under U(1) is the 3-form field strength $G_3 = e^{\Phi/2}(F_3 - \tau H_3)$, which carries charge 1 under U(1).
When computing higher-derivative corrections, one wishes to keep track of the U(1) charge for possible U(1)-violating corrections. In that context it is convenient to work with the complexified vielbein $P$ and 3-form field strength $G_3$.
It is important to distinguish this U(1) from the compensating U(1) used to remove the redundant scalar $\theta$ in the coset representative $\mathcal{V}$. The compensating U(1) is an \SLZ{} transformation, under which the effective action should always be invariant. For example, the modular form $f(\tau, \bar\tau)$ also carries a U(1)-like charge under \SLZ{}, but it does not carry local U(1) charge, and hence may contribute to \SLZ{}-invariant but U(1)-violating corrections.
